\documentclass[11pt]{article}
\usepackage[utf8]{inputenc}
\usepackage[T1]{fontenc}
\usepackage[italian]{babel}
\usepackage{geometry}
\geometry{a4paper, margin=1in}
\usepackage{xcolor}
\usepackage{listings}
\usepackage{hyperref}

\usepackage{tcolorbox}
\tcbuselibrary{breakable}
\begin{document}

\begin{tcolorbox}[colback=blue!5!white,colframe=blue!75!black,title=\textbf{DOMANDA}, breakable]
Esercizio 5\\
Modificare il sorgente del codice per fare in modo che ad ogni utente connesso alla chat arrivi nella console il messaggio ``l'utente sta scrivendo...'' .

NOTA: Lato client, bisogna spedire al server un evento apposito (ad es. ``typing'') quando l'utente scrive sulla tastiera (catturando l'evento di sistema ``keydown''). Lato server, la chiamata \texttt{webSocket.broadcast.emit('typing', data)} rilancia l'evento ``typing'' a tutti i client connessi tranne che a quello dalla quale si è ricevuto il messaggio. Lato client occorre infine gestire la ricezione del messaggio ``typing'' che arriva dal server (vedere gestione del messaggio ``UploadChat'').
\end{tcolorbox}

\begin{tcolorbox}[colback=green!5!white,colframe=green!75!black,title=\textbf{RISPOSTA}, breakable]
\textbf{Esercizio 38 - Segnalazione "l'utente sta scrivendo..."}

\textbf{Modifiche effettuate al codice sorgente:}

\begin{enumerate}
    \item \textbf{Lato Client (\texttt{public/chat.js}):}
    \begin{itemize}
        \item È stato aggiunto un listener (\texttt{'keydown'}) sull'input testuale per intercettare la pressione dei tasti. Quando l'utente digita, viene emesso un evento \texttt{'typing'} inviando al server il nome del mittente:
\begin{lstlisting}[language=JavaScript, basicstyle=\ttfamily\small]
message.addEventListener('keydown', function(){
    webSocket.emit('typing', sender.value);
});
\end{lstlisting}
        \item È stato aggiunto un listener per gestire l'evento \texttt{'typing'} in arrivo dal server. Per evitare che il messaggio sovrascriva eventuali altri utenti che stanno digitando contemporaneamente, viene verificato se la stringa è già presente nel div "feedback" e, se non lo è, viene aggiunta (accodata) e viene scritto un log nella console:
\begin{lstlisting}[language=JavaScript, basicstyle=\ttfamily\small]
webSocket.on('typing', function(data){
    if (feedback.innerHTML.indexOf(data + ' sta scrivendo') === -1) {
        feedback.innerHTML += '<p><em>' + data + ' sta scrivendo...</em></p>';
        console.log(data + " sta scrivendo...");
    }
});
\end{lstlisting}
        \textbf{(Nota:} L'evento \texttt{'UploadChat'} preesistente pulisce già il div "feedback" azzerandone il contenuto quando un messaggio viene effettivamente inviato, rimuovendo quindi la scritta "sta scrivendo...").
    \end{itemize}

    \item \textbf{Lato Server (\texttt{server.js}):}
    \begin{itemize}
        \item È stata aggiunta la gestione dell'evento \texttt{'typing'}. Quando il server lo riceve, usa la funzione "broadcast" per re-inviare lo stesso evento a tutti gli altri client connessi, escludendo il client mittente:
\begin{lstlisting}[language=JavaScript, basicstyle=\ttfamily\small]
webSocket.on('typing', function(data){
    webSocket.broadcast.emit('typing', data);
});
\end{lstlisting}
    \end{itemize}
\end{enumerate}
\end{tcolorbox}

\end{document}
